\subsubsection{Description}
Cellular neural networks are able to make logical operations on binary pictures, this example shows logical negation. Black pixels to white and white to black. It uses a binary picture as an input, matrix setting shown bellow and fixed value boundary condition to produce a binary picture with the edges.
\subsubsection{Setup}

\textbf{Input:} Unimportant (all zeros)\\
\textbf{Boundary conditions:} Fixed.\\
\textbf{Initial output:} Grey-scale picture

\begin{minipage}{0.9\linewidth}
\begin{equation}
A =
\begin{bmatrix}
 0 &  0 &  0 \\
  0 &  1 &  0 \\
  0 &  0 &  0
\end{bmatrix}
B =
\begin{bmatrix}
 0 & 0 & 0 \\
 0 & -2 & 0 \\
 0 & 0 & 0
\end{bmatrix}
Z = 0
\end{equation}
\captionof{figure}{Chosen values of A,B and Z for this experiment}
\end{minipage}

\subsubsection{Results}


\begin{minipage}{0.5\linewidth}
	\centering
	\includegraphics[width=0.9\linewidth]{./Experiments/LOGNOT/fig/Input.png} 
	\captionof{figure}{Input}
\end{minipage}
\begin{minipage}{0.5\linewidth}
	\centering
	\includegraphics[width=0.9\linewidth]{./Experiments/LOGNOT/fig/Output.png}
	\captionof{figure}{Output}
\end{minipage}

\subsubsection{Other logical functions}